%% Conclusion section

\section{Conclusion}
\label{sec:conclusion}
We presented \approach{}, a retrieval-augmented generation (RAG) framework for DeFi attack detection that addresses key limitations of existing approaches. \approach{} constructs an incrementally expandable knowledge base of historical attacks with behavioral summaries, employs a two-stage retrieval mechanism with dynamic similarity thresholding, \Blue{applies adaptive prompting strategies, and incorporates XGBoost post-filtering.} Together, these components bridge the gap between low-level execution traces and high-level attack semantics, while mitigating LLM hallucinations via evidence-grounded reasoning.

\Blue{On 728 balanced \dsystem{} transactions, the Kimi-k3 configuration achieves 93.27\% accuracy, 97.87\% precision, 88.46\% recall, and 92.93\% F1. Compared with direct XGBoost classification, it improves F1 by 6.45 percentage points and reduces false positives from 60 to 7. The Kimi ablation shows a 14.99-point F1 difference between the complete configuration and the run without XGBoost. Two-stage retrieval does not independently improve F1 in the no-XGBoost comparison, highlighting interactions between retrieval and decision refinement. The complementary price-manipulation experiment detects 86 of 91 attacks. Together, these results show the value of combining LLM-based semantic analysis with numerical filtering, while the retrieval ablation identifies a precision--recall trade-off that depends on the pipeline configuration.}

\approach{} faces limitations in scenarios involving cryptographic exploits without behavioral signatures, complex multi-step cross-contract attacks with fragmented semantic signals, and truly novel attack patterns that are absent from the knowledge base. Future work includes integrating formal verification for cryptographic exploit detection, developing more sophisticated call-graph analysis for cross-contract reasoning, and implementing automated knowledge base expansion to remain responsive to emerging threats.

\approach{} demonstrates that retrieval-augmented generation is a promising paradigm for security analysis tasks that require both domain knowledge and robust reasoning. As DeFi protocols continue to evolve, approaches that leverage accumulated security knowledge while adapting to new patterns will be essential for protecting the assets at stake in this ecosystem.
